% Part: incompleteness
% Chapter: arithmetization-syntax
% Section: coding-formulas

\documentclass[../../../include/open-logic-section]{subfiles}

\begin{document}

\olfileid{inc}{art}{frm}
\olsection{\printtoken{P}{formula} यांचे सांकेतीकरण}

$\fn{Term}(x)$ हा संबंध आदिम पुनरावर्ती रीतीने परिभाषित केल्यावर
!!{formula}s साठीचा संबंधित संबंध $\fn{Frm}(x)$ हाही आदिम
पुनरावर्ती रीतीने परिभाषित करता येतो.

\begin{prop}
$x$ ही आण्विक !!{formula} ची गोडेल संख्या असेल, तर आणि तरच
खरे असणारा $\fn{Atom}(x)$ हा संबंध आदिम पुनरावर्ती आहे.
\end{prop}

\begin{proof}
$x$ ही आण्विक !!{formula} ची गोडेल संख्या असेल, तर आणि तरच
पुढीलपैकी एक अट खरी असते:
\begin{enumerate}
\item $n$, $j < x$ आणि $z < x$ अशी मूल्ये अस्तित्वात आहेत की
  प्रत्येक $i < n$ साठी $\fn{Term}((z)_i)$ खरे आहे आणि $x = $
\[
\Gn{\Obj P^n_j(} \concat \fn{flatten}(z) \concat \Gn{)}.
\]
\item $z_1, z_2 < x$ अशी मूल्ये अस्तित्वात आहेत की $\fn{Term}(z_1)$,
  $\fn{Term}(z_2)$ खरे आहेत आणि $x = $
\[
\Gn{{\eq}(} \concat z_1 \concat \Gn{,} \concat z_2 \concat{\Gn{)}}.
\]
  \tagitem{prvFalse}{$x = \Gn{\lfalse}$.}{}
  \tagitem{prvTrue}{$x = \Gn{\ltrue}$.}{}
\end{enumerate}
\end{proof}

\begin{prop}
\ollabel{prop:frm-primrec}
$x$ ही !!a{formula} ची गोडेल संख्या असेल, तर आणि तरच
खरे असणारा $\fn{Frm}(x)$ हा संबंध आदिम पुनरावर्ती आहे.
\end{prop}

\begin{proof}
चिन्हांचा अनुक्रम~$s$ हे !!a{formula} असेल, तर आणि तरच
!!{formula} ची अशी रचनाक्रमिका~$s_0$, \dots, $s_{k-1} = s$
अस्तित्वात असते की जिच्यात रचनानियमांनुसार आण्विक !!{formula}s
पासून~$s$ कसे बनले याची नोंद असते. प्रत्येक $s_i$ चा संकेतांक
$s$ च्या संकेतांक~$x$ पेक्षा मोठा नसतो; पण संपूर्ण रचनाक्रमिका
$\tuple{s_0, \dots, s_{k-1}}$ हिचा संकेतांक त्यापेक्षा लहान
असणे आवश्यक नाही: मूळ इंग्रजीतील उलट दावा एका आण्विक सूत्रासाठीही
खोटा ठरतो. त्याऐवजी मागील पदांच्या सिद्धतेप्रमाणे त्याला
त्या अंतिम सूत्राच्या संकेतांकावर अवलंबून आदिम पुनरावर्ती वरची
मर्यादा देता येते.
\end{proof}

\begin{prob}
\olref[inc][art][frm]{prop:frm-primrec} ची, 
\olref[inc][art][trm]{prop:term-primrec} च्या पहिल्या सिद्धतेच्या
धर्तीवर, सविस्तर सिद्धता द्या.
\end{prob}

\begin{prop}
\ollabel{prop:freeocc-primrec}
$x$ ही गोडेल संख्या असलेल्या सूत्रातील $i$-वे चिन्ह हे
गोडेल संख्या~$z$ असलेल्या चराची मुक्त आस्थिती असेल,
तर आणि तरच खरे असणारा $\fn{FreeOcc}(x, z, i)$ हा संबंध
आदिम पुनरावर्ती आहे.
\end{prop}

\begin{proof}
सराव.
\end{proof}

\begin{prob}
\olref[inc][art][frm]{prop:freeocc-primrec} सिद्ध करा.
!!a{formula} च्या ज्या उपचिन्हमाला स्वतः !!a{formula} आहेत,
त्या त्याची उप-!!{formula} आहेत, हे तथ्य वापरू शकता.
\end{prob}

\begin{prop}
$x$ ही !!a{sentence} ची गोडेल संख्या असेल, तर आणि तरच
खरे असणारा $\fn{Sent}(x)$ हा गुणधर्म आदिम पुनरावर्ती आहे.
\end{prop}

\begin{proof}
!!{sentence} म्हणजे !!{variable}s च्या मुक्त आस्थिती नसलेले
!!a{formula}. म्हणून $\fn{Sent}(x)$ तेव्हाच खरे असते जेव्हा
\begin{multline*}
\bforall{i<\len{x}}{\bforall{z<x}{}}\\
(\bexists{j<z}{z=\Gn{\Obj v_j}} \lif \lnot\fn{FreeOcc}(x,z,i)).
\end{multline*}
\end{proof}


\end{document}
